\chapter*{Abstract}

Electroencephalography (EEG) is a cornerstone tool in clinical and experimental neuroscience, yet it frequently suffers from channel loss due to contact failures, motion artifacts, or amplifier saturation. Interpolating these missing channels is a critical preprocessing step, traditionally addressed through geometric methods such as spherical splines, implemented in reference libraries like MNE-Python. However, these methods operate exclusively in the spatial domain, ignoring the brain's true functional connectivity and the temporal dynamics of the signals.

This thesis presents a reproducible and exhaustive benchmark comparing eighteen EEG channel interpolation methods —geometric, statistical, and Graph Signal Processing (GSP)-based— under a rigorous experimental protocol. The design encompasses three heterogeneous databases (64-channel PhysioNet, 22-channel BCI Competition IV 2a, and 60-channel MNE Sample), 120 loss scenarios combining random and contiguous modes at severities from 10\% to 40\%, and a battery of six metrics covering amplitude error (MAE, RMSE, SNR), morphological fidelity (DTW), and spectral preservation (LSD, Coherence). To ensure a fair comparison, all methods —including classical baselines— were optimized via Optuna with over 15,000 parametric configurations explored, and a confirmatory validation against MNE-Python's community reference implementation was incorporated.

The results support a conditional conclusion. In the balanced confirmatory evaluation, fixed TRSS improves median MAE by 12.4\% over MNE-Python, reaches a 72.0\% win rate, and improves NRMSE by 13.9\%. The advantage grows under grouped or peripheral severe losses, where temporal continuity compensates for missing local spatial information. Nevertheless, MNE-Python retains important advantages: it achieves better global LSD, is substantially faster, and outperforms TRSS on spatially smooth signals. We conclude that TRSS is preferable when accurate reconstruction under non-trivial channel loss is the priority, whereas MNE remains the robust choice for standard preprocessing, LSD-sensitive spectral analysis, and time-constrained pipelines.

All code, experimental configurations, optimization logs, and result artifacts are preserved in the project repository, ensuring full traceability of the findings.

\vspace{1em}
\noindent\textbf{Keywords:} EEG, missing channels, interpolation, Graph Signal Processing (GSP), temporal regularization, TRSS, spherical splines, MNE-Python, Bayesian optimization, Optuna, reproducible benchmark.